\usepackage{xcolor}
\usepackage[standard,thmmarks,amsmath]{ntheorem}
\newcounter{thequestion}
\setcounter{thequestion}{0}
\newcounter{thesubquestion}
\newcounter{thesubsubquestion}
\newtheoremstyle{problem}%
{}{}%
%{\itshape}{}%
%{\bfseries}{}%
%{\newline}{}
%{ \setcounter{thesubquestion}{0}}

\theoremstyle{break}
\theorembodyfont{\rm}
\theoremprework{\vskip 1cm \hrule
                \setcounter{thesubquestion}{0}}
\theorempostwork{\hrule}
\newtheorem{question}{Problem}

\theoremstyle{plain}


%{ 
%%			\refstepcounter{thequestion}
%%           \stepcounter{thequestion}
%           \setcounter{thesubquestion}{0}
%           }
%           {
%           \vskip 0.1cm
%           \vskip 0.3cm
%           }

\newcommand{\subquestion}{\vskip 0.25cm
                \stepcounter{thesubquestion}
                \noindent(\arabic{thesubquestion}) 
                \setcounter{thesubsubquestion}{0}
         }

\newcommand{\subsubquestion}{\vskip 0.10cm
\stepcounter{thesubsubquestion} 
%\noindent(\arabic{thequestion}.\arabic{thesubquestion}.\alph{thesubsubquestion})
\par(\alph{thesubsubquestion})
         }


\definecolor{myred}{rgb}{.5,0,0}


\newif\ifsolutions
\newcommand{\solution}[1]{{
\ifsolutions
\color{myred}
%\small
\vskip 0.75cm
                \setcounter{thesubquestion}{0}
                \noindent {\bf Solution}:{#1}
\fi
}}

\newif\ifuselabel
\newcommand{\mylabel}[1]{{
\ifuselabel
\label{#1}
\fi
}}
